\section{The full-local cone and finite certificates for frame cohomology}
\label{sec:cubic-frame-cohomology-certificate}

The coefficient prime is an odd prime \(p\), and \(k\) is a finite field of characteristic \(p\). All cohomology is continuous. A strict class is zero in \(H^1\) at every place above \(p\), and is unramified elsewhere after its finite coefficient representation has been trivialized. This kernel must be distinguished from the source's cone group.

\begin{proposition}[AC-G76-duality: exact class sectors and the full-local term]
\label{prop:ACG76duality}
Let \(F\) be totally real. Let \(\rho_0\) be a finite representation of \(G_F\) whose image has order prime to \(p\). Let \(N\) be the closed normal subgroup generated by
\[
 I_v\cap\ker\rho_0,\qquad v\nmid p,
\]
and put \(\Gamma=G_F/N\). Let \(V\) be a finite \(k\)-representation factoring through \(\rho_0\). Choose a finite Galois extension \(L/F\), of degree prime to \(p\), trivializing \(V\), and put
\[
 \Delta=\operatorname{Gal}(L/F),\qquad
 A_L=\operatorname{Cl}(\mathcal O_L[1/p])_{(p)}/p.
\]
The strict subgroup of \(H^1(\Gamma,V)\) is canonically
\[
 \operatorname{Hom}_{k[\Delta]}(A_L\otimes_{\mathbf F_p}k,V).
 \tag{G76.1}
\]
One may test strictness at a finite extension of each \(F_v\) of degree prime to \(p\), and may repeat these tests, without changing this subgroup.

For a character \(\xi\), denote the dimension of (G76.1) by \(a_F(\xi)\). If all source local \(H^0(k(\xi))\) vanish, the cone \(H^1_C(\mathrm{loc},k(\xi),0)\) has dimension \(a_F(\xi)\). For the full local condition and trivial coefficients,
\[
 \begin{split}
 H^1_C(\mathrm{loc},k,\mathrm{top})&=H^1(G_{F,S_p},k),\\
 \dim_k H^1(G_{F,S_p},k)&=1+a_F(\bar\kappa)+u_F,\\
 u_F&=\sum_{v\mid p}\dim_k k(1)^{G_{F_v}}
       -\dim_k k(1)^{G_F}\\
    &=\#\{v\mid p:\mu_p\subset F_v\}.
 \end{split}
 \tag{G76.2}
\]
Here \(a_F(\bar\kappa)\) can be computed with \(L=F(\mu_p)\), or a larger extension of degree prime to \(p\). The groups defining \(u_F\) are the full \(G_{F_v}\), even when a frame presents strict conditions using finite local extensions.
\end{proposition}

\begin{proof}
Inflation--restriction identifies
\(H^1(G_F,V)=H^1(G_L,V)^\Delta\), since the positive-degree cohomology of \(\Delta\) vanishes in characteristic \(p\). On \(G_L\), a cocycle is a homomorphism. Restricting a cocycle which factors through \(\Gamma\) kills every inertia group outside \(p\): its restriction to inertia in \(\ker\rho_0\) vanishes, and the remaining quotient has order prime to \(p\). Strictness at \(p\) says that the homomorphism kills all decomposition groups above \(p\). Global class field theory gives (G76.1). Its equivariance is
\(f(\delta h\delta^{-1})=\delta f(h)\), fixing the character orientation.

Conversely, a homomorphism in (G76.1) descends to a class of \(G_F\). On \(\ker\rho_0\), its representing cocycle is a homomorphism. On inertia there it vanishes after restriction to \(G_L\); the remaining quotient has order prime to \(p\), so it vanishes already. It therefore kills the generators of \(N\) and their normal closure, and factors through \(\Gamma\). Restriction to a local extension of degree prime to \(p\) is injective on \(H^1\), by restriction--corestriction. This proves all statements about local tests.

The cone sequence of AC-G74-cone is
\[
 0\longrightarrow
 \operatorname{coker}\bigl(H^0(F,V)\longrightarrow
                 \textstyle\bigoplus_{\mathrm{loc}}H^0(V)\bigr)
 \longrightarrow H^1_C
 \longrightarrow H^1_{\mathrm{strict}}\longrightarrow0.
 \tag{G76.3}
\]
It gives the strict-cone assertion. For the condition \(\mathrm{top}\), every auxiliary local tuple disappears modulo cone boundaries, leaving \(H^1(\Gamma,k)\). An additive \(k\)-character of \(\Gamma\) kills all away-\(p\) inertia, because its allowed residual inertia quotient has order prime to \(p\). Conversely every character unramified away from \(p\) kills \(N\). This proves the first equality of (G76.2).

Tate's global Euler formula for trivial coefficients over totally real \(F\) gives
\[
 1-\dim H^1(G_{F,S_p},k)+\dim H^2(G_{F,S_p},k)=0.
 \tag{G76.4}
\]
Identify the finite dual of \(k\) with \(k\) by the trace pairing. Poitou--Tate and local Tate duality give
\[
 \begin{split}
 \dim H^2(G_{F,S_p},k)
 &=\dim \operatorname{Sha}^1(G_{F,S_p},k(1))\\
 &\quad+\sum_{v\mid p}\dim H^0(F_v,k(1))
                   -\dim H^0(F,k(1)).
 \end{split}
 \tag{G76.5}
\]
Indeed the kernel of localization on \(H^2(k)\) is dual to \(\operatorname{Sha}^1(k(1))\). Its image is the kernel of the surjective final local-to-global map in the Poitou--Tate sequence, whose dimension is the displayed difference of local and global invariant dimensions. At real places the Tate-modified groups vanish because \(p\) is odd. For this application one may adjoin \(\bar\kappa\) to \(\rho_0\): it is unramified away from \(p\), so the defining subgroup \(N\) is unchanged. Apply (G76.1) to \(k(1)\). The totally real field \(F\) contains no primitive \(p\)-th root of unity, and each local invariant space has dimension zero or one. This proves (G76.2).
\end{proof}

\begin{proposition}[AC-G76-host: bounded full-local cohomology on the dihedral host]
\label{prop:ACG76host}
Let \(F_n\) be the layers of the cyclotomic \(\mathbf Z_p\)-extension of a totally real field \(F\). Suppose either \(F/\mathbf Q\) is abelian, or \(F=LM\) is the real part of AC-G73-host: \(L/\mathbf Q\) is dihedral with cyclic subgroup of order \(p^s\), its fixed quadratic field is \(E\), and \(M/\mathbf Q\) is real multiquadratic. In the second case retain \(p>3\).

There is \(C_h\), independent of \(n\), such that for every quotient \(\Gamma_n\) as above with residual image of order prime to \(p\),
\[
 \dim_k H^1_C(\mathrm{loc}_n,k,\mathrm{top})\le 1+C_h.
 \tag{G76.6}
\]
No finiteness or \(\mu=0\) assertion about the cubic coefficient splitting field is made.
\end{proposition}

\begin{proof}
Put \(Q_F=F(\mu_p)\). In the abelian case \(Q_F/\mathbf Q\) is abelian. In the dihedral case
\[
 Q_F=LM(\mu_p),\qquad Q_0=EM(\mu_p)
\]
have \(Q_0/\mathbf Q\) abelian and \(Q_F/Q_0\) cyclic of degree dividing \(p^s\). Ferrero--Washington gives cyclotomic \(\mu(Q_0)=0\). Iwasawa's propagation theorem for cyclic degree-\(p\) extensions of cyclotomic \(\mathbf Z_p\)-fields, iterated along this extension, gives \(\mu(Q_F)=0\). Infinite-place ramification is absent since \(p\) is odd. Thus the maximal unramified abelian pro-\(p\) group \(X(Q_{F,\infty})\) is finitely generated over \(\mathbf Z_p\).

The precise propagation can also be checked as follows. For finite \(p\)-group \(P=\operatorname{Gal}(Q_{F,\infty}/Q_{0,\infty})\), the maximal unramified abelian pro-\(p\) extension of \(Q_{F,\infty}\) is Galois over \(Q_{0,\infty}\); call the latter group \(H\). Its inertia groups inject into \(P\). Only finitely many primes ramify, since every finite prime is finitely decomposed in a cyclotomic tower. Their images generate a finite subgroup of \(H^{\mathrm{ab}}\), and the quotient is a quotient of \(X(Q_{0,\infty})\). Hence \(H\) has finite Frattini quotient and is finitely generated pro-\(p\). Schreier's theorem makes its open abelian subgroup \(X(Q_{F,\infty})\) finitely generated over \(\mathbf Z_p\). This checks the needed implication without using a \(\lambda\)-formula.

The \(p\)-split quotient \(X'(Q_{F,\infty})\) is also finitely generated over \(\mathbf Z_p\). Set
\[
 D=\dim_{\mathbf F_p}X'(Q_{F,\infty})/p.
\]
Start after all \(p\)-primes are totally ramified. AC-G74's finite-layer comparison gives a surjection from coinvariants of \(X'(Q_{F,\infty})\) onto
\(\operatorname{Cl}(\mathcal O_{F_n(\mu_p)}[1/p])_{(p)}\).
Thus its \(\bar\kappa\)-sector modulo \(p\) has dimension at most \(D\). The number \(s_p\) of \(p\)-primes of \(F_n\) is constant thereafter. Formula (G76.2) gives (G76.6) with \(C_h=D+s_p\); enlarge this for the finitely many earlier layers. The bound depends on the underlying field, not on a single-power residual seed.

For the AC-G73-host tower at \(p=269\), every finite layer has degree of the form \(2^a269^b\). If a completion contained \(\mu_{269}\), its degree over \(\mathbf Q_{269}\) would be divisible by \(268=4\cdot67\), whereas that local degree divides the global degree. This is impossible. Consequently \(u_{F_n}=0\) at every layer, and after the initial layer chosen above one may take \(C_h=D\).
\end{proof}

\begin{corollary}[AC-G76: the remaining finite cohomological certificates]
\label{cor:ACG76}
Suppose the source cohomological frame clauses have their stated arithmetic realization over a totally real \(F\), with CM quadratic field \(B/F\), character \(\delta\), and character \(\Psi\). Retain exactly the \(\mathrm{GalLink}\) kernel and local-place hypotheses and \(\mathrm{Frame.Ok.f3}\). This is also a criterion for an explicitly proved replacement with the same cohomological clauses; no such replacement is inferred here.

Choose a common extension of degree prime to \(p\) splitting \(\Psi,\delta,\bar\kappa\). The five cohomological clauses are precisely
\[
\begin{array}{c|c}
 \text{source clause}&\text{class-sector condition}\\ \hline
 \mathrm{f4}&a_F(\Psi^{-1}\delta)>0\\
 \mathrm{f4'}&a_F(\Psi^{-2}\delta)>0\\
 \mathrm{b1}&a_F(\Psi^{-1})+a_F(\Psi^{-1}\delta)\le C_0\\
 \mathrm{b1p}&a_F(\Psi^{-2}\delta)\le C_0\\
 \mathrm{h1}&a_F(\bar\kappa)+u_F\le C_0 .
\end{array}
\tag{G76.7}
\]
In particular \(\mathrm{f4'}\) and \(\mathrm{b1p}\) concern the same finite sector.

Start the cyclotomic tower after total ramification at all \(p\)-primes. If the first two sectors are positive initially, they remain positive. Suppose AC-G74 supplies
\[
 \begin{split}
 a_{F_n}(\Psi^{-1})+a_{F_n}(\Psi^{-1}\delta)&\le R_1p^n+c_1,\\
 a_{F_n}(\Psi^{-2}\delta)&\le R_2p^n+c_2,
 \end{split}
 \qquad d_n=d_0p^n,
 \tag{G76.8}
\]
with constants independent of \(n\). If AC-G76-host applies and
\[
 d_0>3\max(R_1,R_2),
 \tag{G76.9}
\]
then these five clauses and both numerical conditions of \(\mathrm{Frame.Big}\) hold at every sufficiently large layer for a suitable \(C_{0,n}\). The \(R_i\) may be supplied by AC-G74's finite class-sector plus permutation-augmentation certificates. All other \(\mathrm{Frame.Ok}\) fields, the seed and host requirements, and the actual replacement of one character and its powers by two cubic adjacent coefficients remain separate inputs.
\end{corollary}

\begin{proof}
At each local map \(\mathrm{f3}\) makes \(\Psi^2\) nontrivial, and the map lands in \(\ker\delta\). Thus the local invariants of \(\Psi^{-1}\), \(\Psi^{-1}\delta\), and \(\Psi^{-2}\delta\) vanish. Equation (G76.3) has no correction, even with repeated maps. These maps cover the \(p\)-places up to extensions of degree at most two. Equation (G76.1) proves the first, second, and fourth rows.

For the third row \(\mathrm{loc2}\) gives both \(p\)-embeddings into the CM quadratic field. At an inert place two tests may be conjugate; this changes neither the strict kernel nor its cone, since local \(H^0\) is zero. Quadratic induction gives
\[
 \operatorname{Ind}_{G_B}^{G_F}k(\Psi^{-1}|_{G_B})
       =k(\Psi^{-1})\oplus k(\Psi^{-1}\delta).
\]
It is compatible with all strict local tests, because \(2\) is invertible. This proves the sum. The last row is (G76.2).

A nontrivial character of order prime to \(p\) retains its image on a subgroup of \(p\)-power index. Its invariant space there remains zero. Inflation--restriction therefore makes restriction on its \(H^1\) injective at each cyclotomic step. Strictness is preserved by restriction, so the two initial positive certificates persist. The same order argument preserves \(\mathrm{f3}\).

Put \(R=\max(R_1,R_2)\), choose an integer \(C\) dominating \(c_1,c_2,C_h,0\), and put \(C_{0,n}=Rp^n+C\), increasing the constant for rounding if necessary. Equations (G76.6)--(G76.8) and positivity prove the five clauses. Equation (G76.9) implies \(3C_{0,n}+7<d_0p^n\) eventually. Total ramification multiplies every positive local degree by \(p^n\), so all \(d_{v,n}>5\) eventually. The other frame fields have not been used or supplied.
\end{proof}

\paragraph{Verified source locators and scope.}
The pinned \path{FrobeniusSmoothVectors/ContinuousCoboundaries2.lean} gives
\(\mathrm{H1C}\), \(\mathrm{Frame.loc2}\), \(\mathrm{Frame.Ok}\)
(\(\mathrm{f3,f4,f4',b1,b1p,h1}\)), and \(\mathrm{Frame.Big}\).
The kernel and local realization are
\path{CurveDVRSpecialization/HostFrameEx.lean},
\(\mathrm{GalLink.p\_ker}\), its character compatibilities and local-place fields.
The pinned commit is \texttt{7cbf7a0c32fca1badf05931879a37c3f8a2b03b4}.
Milne, \emph{Arithmetic Duality Theorems}, second edition, I, Theorem 5.1,
printed p.\ 67, and Remark 5.2(c), printed p.\ 68, give the Euler formula
and its precise duality consequence. Hida's \S3.5, Theorem 3.12 has the
same correct statement, but its printed general Euler-duality calculation
contains an apparent error; it is not used as a verified proof.
Sharifi, \emph{Iwasawa Theory}, Appendix A, Theorems A.0.3, A.0.6 and A.0.8,
gives the finite duality statements. His Theorem 6.2.1 states
Ferrero--Washington; Proposition 3.1.11, including its printed proof,
gives finite decomposition and eventual total ramification.
Schettler, \emph{Generalizations of Iwasawa's Riemann--Hurwitz Formula for
Cyclic \(p\)-Extensions of Number Fields}, Theorem 2, printed pp.\ 1--2,
states the propagation theorem and cites Iwasawa,
\emph{T\^ohoku Math.\ J.} 33 (1981), 263--288.
The original Iwasawa PDF was inaccessible in this check; no internal
locator or fresh reading of its proof is claimed.

\paragraph{Geometric dictionary and retained gate.}
The class terms classify abelian finite etale \(p\)-covers splitting at the
specified \(p\)-adic places. Their coefficient sectors and restriction maps
use the already-audited SGA~1, Expose~V, \S\S4 and 7 framework; no new
archive reading is claimed. AC-G76-host removes the full-local growth
condition on the dihedral host. Positive cubic class certificates, the
two actual adjacent-root upper inequalities, and a valid arithmetic
replacement for the literal frame remain unsupplied.
